%!TEX root = ../main.tex

\chapter{Agents: Sufficiency Check and Tool Use}\label{ch:agents}

Passage retrieval misses the gold section for about 27\% of the validation questions, and \fone{}
there is far lower. The graph facts rescue some of these questions (4.6 points on 342 of them,
Chapter~\ref{ch:graphhelp}), but a router could not tell before generation which ones, and
the generator with the dialogue history refuses too rarely, which is why the decision to
abstain comes from the prompt without history (Chapter~\ref{ch:convgen}). This chapter asks
whether the model itself can judge, after reading the retrieved context, that it is not enough,
and whether looking further repairs the answer at an acceptable cost. The baseline is E with
citations: passages with up to 30 facts, the conversational generator
(\texttt{gpt-5.4-nano}, G2 + abstain(G0)) and the citation prompt of Chapter~\ref{ch:sources}.

The experiments were planned in advance, with the stopping rule fixed before any run: a variant is
kept only if the paired interval of the \fone{} difference to E on all questions is above 0 (and
answerable questions do not get significantly worse), or if it ties E ($\Delta \ge -0.2$) at
least 10\% cheaper; it is dropped if the interval on all or on answerable questions lies entirely
below 0. Development and the decision were both done on dev-500 here, because the first runs
already ruled the variants out (see the end of the chapter); no variant was run on the full validation
or the test split.


% ===========================================================================
\section{Variants}\label{sec:agents:variants}

\begin{description}
    \item[A1, fixed controller.] One call checks the rewritten query, the top five sections and
        the facts, and answers in JSON whether the context is sufficient, what is missing and which
        query would search for it. If it is sufficient, the answer is the one of E. If not, there is
        one extra round, never more: the new query goes through the same hybrid retrieval and
        reranking, the five best new sections are added to the old five (duplicates removed, at
        most ten), up to 15 facts of the entities linked only by the new query are added, and the
        question is answered again. The refusal is decided by G0 on the new context; a paired
        alternative refuses if a second check on the new context is still insufficient.
    \item[A2, tool-calling agent.] The model starts from the context of E and has four tools:
        \texttt{search\_passages(query)}, \texttt{get\_entity\_facts(entity)},
        \texttt{answer(text, citations)} and \texttt{refuse()}. At most three calls of the two
        search tools are allowed, after which it must answer or refuse; the loop is plain Python
        over the existing retrieval modules. The refusal is either its own (\texttt{refuse()}) or,
        in addition, the decision of G0 on the final context.
\end{description}
A2 was to be run only if the check of A1 was better than chance on dev-500, defined as Youden's $J
\ge 0.2$ with the lower bound of its bootstrap interval above 0, measured against the label ``a gold
section is in the context'' (among the passages or the evidence section of a fact). Every run
was estimated in calls and cost before it started.


% ===========================================================================
\section{The Sufficiency Check}\label{sec:agents:check}

The check passes the gate: $J = 0.33$ with a 95\% interval of $[0.23, 0.42]$ on the 469 answerable
dev-500 questions (strict label, passages only: $0.29$, $[0.19, 0.37]$). Table~\ref{tab:agents:check}
gives the confusion matrix. When it says the context is sufficient it is right for 86\% of the
questions, but it finds only 63\% of the questions whose context holds the gold section; when it
says insufficient it is right for 39\% (the base rate of a missing gold section is 25\%) and finds
70\% of the missing ones. On the 31 unanswerable questions it says insufficient 58\% of the time.
Two replies of 500 could not be parsed.

\begin{table}[ht]
    \centering
    \caption{The sufficiency check against ``gold section in context'', answerable dev-500
    questions ($n = 469$).}
    \label{tab:agents:check}
    \small
    \begin{tabular}{lrr}
        \toprule
        & Gold in context & Gold missing \\
        \midrule
        Check says sufficient   & 220 &  36 \\
        Check says insufficient & 130 &  83 \\
        \bottomrule
    \end{tabular}
\end{table}


% ===========================================================================
\section{Results on dev-500}\label{sec:agents:results}

\begin{table}[ht]
    \centering
    \caption{Agents on dev-500 (\%, $n = 500$), against E with citations. $\Delta$: paired
    bootstrap difference to E with its 95\% interval. Cost: USD per 1{,}000 turns at the list
    price of \texttt{gpt-5.4-nano}, rewrite left out; calls: model calls per turn (E: 1.9). A2 tokens
    are those reported by the API (tool definitions included), the others are counted from the
    prompts.}
    \label{tab:agents:main}
    \small
    \begin{tabular}{lrrrrrr}
        \toprule
        System & \fone{} & $\Delta$ all & $\Delta$ answerable & Abstain rec. / false & Cost & Calls \\
        \midrule
        E                          & 31.2 & ---                     & ---                      & 41.9 / 4.7  & 1.39 & 1.9 \\
        A1, abstain(G0)            & 30.6 & $-0.6$ ($-1.6$, $0.4$)  & $+0.1$ ($-0.5$, $0.6$)   & 32.3 / 4.5  & 2.29 & 2.9 \\
        A1, own refusal            & 24.6 & $-6.6$ ($-8.5$, $-4.7$) & $-8.5$ ($-10.1$, $-6.9$) & 64.5 / 37.1 & 2.33 & 3.0 \\
        A2, abstain(G0)            & 29.7 & $-1.5$ ($-2.9$, $-0.1$) & $-0.9$ ($-2.0$, $0.1$)   & 32.3 / 4.5  & 1.98 & 2.5 \\
        A2, own refusal            & 28.3 & $-2.8$ ($-4.5$, $-1.3$) & $-0.5$ ($-1.6$, $0.6$)   & 3.2 / 1.3   & 1.23 & 1.5 \\
        \bottomrule
    \end{tabular}
\end{table}

\paragraph{A1.} The extra round was taken by 231 of the 500 questions (46\%). It adds little: on the
213 answerable ones, the new query alone puts a gold section in its top five for 45.5\% of them,
less than the query of E (53.5\%), and the union with the old top five raises the recall of the
passages from 53.5\% to 59.2\% (+5.6 points). Counting the facts as well, a gold section is in the
context for 61.0\% before and 65.7\% after. The larger context makes G0 abstain less (abstain recall
32.3\% against 41.9\%), so the difference to E is $-0.6$ on all questions and $+0.1$ on answerable
ones, at 65\% more cost. Refusing when a second check is still insufficient is much worse: the
second check calls 80\% of the questions of the round insufficient, and false abstention reaches
37\%. With the interval of the difference on all questions at most 0.4 above 0, A1 cannot satisfy
the keep condition at a higher cost than E, and it was not run on the full validation split.

\paragraph{A2.} The agent used a tool on 44\% of the questions (one call for 209, two for 11, none
reached the limit of three); 220 calls were passage searches and 11 fact lookups, and
\texttt{refuse()} was called four times. Tool use is a weaker predictor of a missing gold section
than the check of A1 ($J = 0.21$, $[0.12, 0.31]$). With the abstention decision of G0 the difference
to E is $-1.5$ on all questions ($[-2.9, -0.1]$), a significant loss. With its own refusal the agent
ties E on answerable questions ($-0.5$, $[-1.6, 0.6]$) at 12\% lower cost, since one call replaces
the G0 and G2 pair when no tool is used, but it almost never refuses (abstain recall 3.2\% against
41.9\%), which gives $-2.8$ on all questions. Both variants have the interval on all questions below
0 and were dropped; the 31 unanswerable questions of dev-500 make the abstention figures noisy.


% ===========================================================================
\section{An Exploratory Check}\label{sec:agents:explore}

After the A1 result, one check that needs no \gls{llm} call was added, and it was not part of the plan:
for the questions A1 sent to the extra round, how often is a gold section at ranks 6 to 10 of the
original reranked list of E, compared with the sections the new query brought in?
Table~\ref{tab:agents:explore} shows that going five sections deeper in the original list recovers more
gold sections than the model's own query, and that the two are largely complementary. This was
chosen after seeing the result on dev-500, on one subset and without an interval, so it is a
hypothesis for further work and not a result.

\begin{table}[ht]
    \centering
    \caption{Exploratory: share of the answerable questions called insufficient by the check where
    a gold section outside the top five of E is among the added sections (\%).}
    \label{tab:agents:explore}
    \small
    \begin{tabular}{lrrrr}
        \toprule
        Group & $n$ & Ranks 6--10 of E & New query of A1 & Both \\
        \midrule
        All answerable          & 213 &  8.9 & 5.6 & 12.7 \\
        Gold section in context & 130 &  5.4 & 3.1 &  6.9 \\
        Gold section missing    &  83 & 14.5 & 9.6 & 21.7 \\
        \bottomrule
    \end{tabular}
\end{table}


% ===========================================================================
\section{Conclusion}\label{sec:agents:conclusion}

The model can tell when the context lacks the gold section, but neither a second search with its
own query nor a tool-using agent recovers it often enough to pay off, and the larger context
costs abstentions on the unanswerable questions. An agent that decides to refuse on its own almost
never refuses, as with the dialogue history in G2: the decision to abstain has to come from a
separate call such as G0. E stays the final system. A direction that follows from
Section~\ref{sec:agents:explore} is a second round that adds ranks 6 to 10 of the original list and
keeps a separate refusal decision; it would need a fresh run on the validation split to mean anything.
